\begin{figure*}[t]
  \centering
  \includegraphics[width=\textwidth]{figures/overview-standalone.pdf}
  \caption{\textbf{WitnessRAG architecture.}
  (1) \emph{Memory construction}: an extractor turns dialogue histories into propositions with triples,
  event dates, and source links. (2) \emph{Memory retrieval}: hybrid retrieval supplies evidence to
  a planner, which specifies conjunctive requirements, answer types, a
  temporal reference, and scoring weights. The executor joins graph facts
  under consistent bindings to produce witnesses. A conditional verifier
  checks candidates against source utterances; gaps and rejections feed
  bounded replanning. Within retrieval, accepted witnesses and candidate facts receive
  selection bonuses, while relevance ranking supplies additional evidence.
  The full profile delivers dated propositions and cached passage summaries
  to the reader. Blue outlines denote LLM operations, green outlines denote
  retrieval or deterministic processing, dotted links denote provenance,
  and dashed arrows denote feedback. Witness acceptance does not guarantee
  that all supporting facts survive the final context budget.}
  \Description{The architecture has two stages: memory construction and memory retrieval. Memory construction
  extracts propositions and indexes their time and importance metadata.
  Source passages and the dated proposition graph are connected by
  provenance. A question triggers evidence retrieval and an LLM planner.
  The planner specifies conjunctive requirements, answer type, temporal
  reference and weights. A graph executor finds consistent bindings without
  LLM calls. A conditional verifier checks source utterances, returning
  gaps or rejections to the planner. Accepted support and relevance-ranked
  candidates guide context selection. The reader receives dated statements
  and cached summaries and generates an answer.}
  \label{fig:overview}
\end{figure*}
